The bound $S^*\ge\max\{\hat R-b,0\}$ is the claim we rest on. The result for Cost-Deterred Hiding and Partial Pooling holds only where the market is outside Complete Adverse Selection, and which regime a market is in depends on the hiding cost. With a single cost per market, calibrated to the observed hiding rate, every market is in Partial Pooling at the estimated attention by construction (\attScBaseNPpApt{} of \attNMarketsApt{} apartment markets and \attScBaseNPpHouse{} of \attNMarketsHouse{} house markets), which follows from the calibration; the registry shows interior hiding rates in essentially every market, which in this model is possible only in Partial Pooling. Reading the hiding cost instead off each market's premium, $c_0=\hat\Delta_m$, places \attRawNCasApt{} apartment markets and \attRawNCasHouse{} house markets in Complete Adverse Selection, with mean threshold \attRawRbarRatioApt{} and \attRawRbarRatioHouse{} of $b$ (Appendix~\ref{app:attention}); we do not take that reading, because it returns negative costs in \attRawNNegApt{} apartment and \attRawNNegHouse{} house markets. The bound is what survives both readings: whatever the hiding cost, the subsidy required is at least the shortfall of $\hat R$ below the rating premium $\hat b$. Because $\bar R=\hat b$ in the regimes the data support, the minimum subsidy does not depend on the observed hiding rate or on $\hat\Delta$, which is an implication for how subsidy needs vary across markets; testing it needs programme take-up by commune, which we do not have.
